% 贯穿图中A={1}, B={4}。删除S={2}后仍有路径1-3-4；删除S={2,3}后没有从1到4的路径。
\begin{tikzpicture}[>=Stealth,
  v/.style={circle,draw=FigureStroke,fill=FigureBlueFill,line width=.8pt,minimum size=6.5mm,inner sep=0pt},
  va/.style={v,fill=FigureGreenFill,draw=FigureGreen},
  vb/.style={v,fill=FigureYellowFill,draw=FigureYellow},
  vs/.style={circle,draw=FigureRed,dashed,fill=FigureRedFill,line width=.8pt,minimum size=6.5mm,inner sep=0pt},
  e/.style={draw=FigureStroke,line width=.8pt},
  ed/.style={draw=FigureMuted,dashed,line width=.6pt},
  h/.style={draw=FigureBlue,line width=1.6pt}]
\begin{scope}
  \node[va] (1) at (-1.2,0) {$1$}; \node[v] (2) at (0,.85) {$2$};
  \node[v] (3) at (0,-.85) {$3$}; \node[vb] (4) at (1.2,0) {$4$};
  \draw[e] (1)--(2) (1)--(3) (2)--(3) (2)--(4) (3)--(4);
  \node[align=center] at (0,-1.65) {(a) $A=\{1\}$，$B=\{4\}$};
\end{scope}
\begin{scope}[xshift=3.6cm]
  \node[va] (1) at (-1.2,0) {$1$}; \node[vs] (2) at (0,.85) {$2$};
  \node[v] (3) at (0,-.85) {$3$}; \node[vb] (4) at (1.2,0) {$4$};
  \draw[ed] (1)--(2) (2)--(3) (2)--(4);
  \draw[h] (1)--(3)--(4);
  \node[align=center] at (0,-1.65) {(b) 删除$\{2\}$：仍有$1,3,4$};
\end{scope}
\begin{scope}[xshift=7.2cm]
  \node[va] (1) at (-1.2,0) {$1$}; \node[vs] (2) at (0,.85) {$2$};
  \node[vs] (3) at (0,-.85) {$3$}; \node[vb] (4) at (1.2,0) {$4$};
  \draw[ed] (1)--(2) (1)--(3) (2)--(3) (2)--(4) (3)--(4);
  \node[align=center] at (0,-1.65) {(c) 删除$\{2,3\}$：已分隔};
\end{scope}
\end{tikzpicture}
